% !PRINTED 47
\weberpage{47}
%<<CONT

\[
M_1=\Delta_0=f(1)-f(0).
\]
Setzen wir
\begin{equation}
\begin{aligned}
 \Delta_{x+1}-\Delta_x&=\Delta_x',\\
 \Delta_{x+1}'-\Delta_x'&=\Delta_x'',\\
 &\cdots\\
 \Delta_{x+1}^{(n-2)}-\Delta_x^{(n-2)}&=\Delta_x^{(n-1)}.
\end{aligned}
\tag{6}
\end{equation}
so wird also hiernach
\[
 M_0=f(0),\quad M_1=\Delta_0,
 \quad M_2=\Delta_0',\quad\ldots,\quad M_n=\Delta_0^{(n-1)},
\]
und die Formel (3) ergiebt
\begin{equation}
 f(x)=f(0)+\Delta_0B_1^{(x)}+\Delta_0'B_2^{(x)}+\Delta_0''B_3^{(x)}+\cdots+\Delta_0^{(n-1)}B_n^{(x)}.
\tag{7}
\end{equation}
und ebenso kann man die Functionen $\Delta_x,\Delta_x',\Delta_x'',\ldots$ ausdrücken:
\begin{equation}
\begin{aligned}
 \Delta_x&=\Delta_0+\Delta_0'B_1^{(x)}+\Delta_0''B_2^{(x)}+\cdots+\Delta_0^{(n-1)}B_{n-1}^{(x)},\\
 \Delta_x'&=\Delta_0'+\Delta_0''B_1^{(x)}+\cdots+\Delta_0^{(n-1)}B_{n-2}^{(x)}.
\end{aligned}
\tag{8}
\end{equation}
Die $\Delta_x,\Delta_x',\Delta_x'',\ldots,\Delta_x^{(n-1)}$ sind ganze rationale Functionen der Grade $n-1,n-2,\ldots,0$, die letzte von ihnen also constant. Um sie alle darzustellen, braucht man nur die Werthe
\[
 f(0),\quad \Delta_0,\quad \Delta_0',\quad \Delta_0'',\quad\ldots,\quad \Delta_0^{(n-1)},
\]
die man am leichtesten berechnet, wenn man eine Tabelle anlegt, die für den Fall $n=3$ z. B. folgende Form haben würde:
\[
\begin{array}{c|c|c|c}
f(0) &&&\\
f(1)&\Delta_0&&\\
f(2)&\Delta_1&\Delta_0'&\\
f(3)&\Delta_2&\Delta_1'&\Delta_0''
\end{array}
\]
und, wenn die $f(0),f(1),\ldots$ gegeben sind, durch einfache Subtractionen berechnet wird.

\begin{center}
{\large \S.\ 11.}

\vspace{0.5em}
{\large\gesperrt{Arithmetische Reihen höherer Ordnung.}}
\end{center}

Das Vorhergehende ist die Grundlage für die Theorie der arithmetischen Reihen höherer Ordnung.

Wir bezeichnen mit
\begin{equation}
u_0,u_1,u_2,u_3,\ldots
\tag{1}
\end{equation}
irgend eine unendliche Reihe von Grössen und mit
%CONT>>
